%! Matrices and Column Spaces — Unit 1, Lesson 1.3 — Slides (Beamer)
\documentclass[aspectratio=169, 11pt]{beamer}
\usepackage{linalg-beamer}   % provides \burgundyheader{} and \sectionlabel[color]{}
\newcommand{\xx}{\mathbf{x}}
\newcommand{\bb}{\mathbf{b}}

\begin{document}

% TITLE SLIDE (CourseName is not defined in beamer, so the name is literal)
{
\setbeamercolor{background canvas}{bg=burgundy}
\begin{frame}[plain]
  \vspace{2.8cm}\hspace{0.55cm}%
  \begin{minipage}{0.9\textwidth}
    {\color{goldacc}\bfseries\small LESSON 1.3}\\[0.5cm]
    {\color{white}\bfseries\LARGE Matrices and Column Spaces}\\[0.3cm]
    {\color{white}\bfseries\large A matrix is vectors in columns --- what can they build?}\\[1.0cm]
    {\color{blushmid}\small Linear Algebra~$\cdot$~Unit 1: Vectors and Matrices}
  \end{minipage}
\end{frame}
}

% HOOK
\begin{frame}[t, plain]
  \burgundyheader{A matrix is just vectors side by side}
  \sectionlabel{Hook}
  \begin{columns}[T]
    \begin{column}{0.58\textwidth}
      Read $A=\begin{bsmallmatrix} 1 & 3 \\ 2 & 1 \end{bsmallmatrix}$ as two columns
      $\begin{bsmallmatrix} 1\\2 \end{bsmallmatrix}$ and $\begin{bsmallmatrix} 3\\1 \end{bsmallmatrix}$.
      Multiplying by $\xx$ \emph{combines} them:
      \[
        A\begin{bsmallmatrix} 2\\1 \end{bsmallmatrix}
        = 2\begin{bsmallmatrix} 1\\2 \end{bsmallmatrix} + 1\begin{bsmallmatrix} 3\\1 \end{bsmallmatrix}
        = \begin{bsmallmatrix} 5\\5 \end{bsmallmatrix}.
      \]
      \textbf{Question.} As $\xx$ ranges over everything, which targets $\bb$ can we build?
    \end{column}
    \begin{column}{0.40\textwidth}
      \centering
      \begin{tikzpicture}[scale=0.5]
        \draw[step=1, black!15, thin] (-0.3,-0.3) grid (5.3,5.3);
        \draw[->, thick, black!60] (-0.4,0)--(5.4,0);
        \draw[->, thick, black!60] (0,-0.4)--(0,5.4);
        \draw[->, very thick, burgundy] (0,0)--(1,2) node[left]{\scriptsize $\begin{bsmallmatrix} 1\\2 \end{bsmallmatrix}$};
        \draw[->, very thick, royalblue] (0,0)--(3,1) node[right]{\scriptsize $\begin{bsmallmatrix} 3\\1 \end{bsmallmatrix}$};
        \draw[->, very thick, black!70, dashed] (0,0)--(5,5) node[above right]{\scriptsize $\begin{bsmallmatrix} 5\\5 \end{bsmallmatrix}$};
      \end{tikzpicture}
    \end{column}
  \end{columns}
\end{frame}

% MATRIX AS COLUMNS + Ax
\begin{frame}[t, plain]
  \burgundyheader{$A\xx$ combines the columns}
  \sectionlabel{Definition}
  A matrix is vectors lined up as columns; $A\xx$ weights them by $\xx$ and adds:
  \[
    A\xx = \begin{bmatrix} \mathbf{a}_1 & \mathbf{a}_2 \end{bmatrix}\begin{bmatrix} x_1 \\ x_2 \end{bmatrix}
    = x_1\,\mathbf{a}_1 + x_2\,\mathbf{a}_2 .
  \]
  \begin{itemize}
    \item This is a \textbf{linear combination of the columns} --- the Lesson 1.1 idea, repackaged.
    \item Example: $\begin{bsmallmatrix} 4 & 0 \\ 1 & 5 \end{bsmallmatrix}\begin{bsmallmatrix} 1\\2 \end{bsmallmatrix}=1\begin{bsmallmatrix} 4\\1 \end{bsmallmatrix}+2\begin{bsmallmatrix} 0\\5 \end{bsmallmatrix}=\begin{bsmallmatrix} 4\\11 \end{bsmallmatrix}$.
  \end{itemize}
\end{frame}

% COLUMN SPACE
\begin{frame}[t, plain]
  \burgundyheader{The column space: everything the columns build}
  \sectionlabel{Column space}
  \begin{columns}[T]
    \begin{column}{0.55\textwidth}
      The \textbf{column space} $C(A)$ is the set of \emph{all} $A\xx$ --- every combination of the
      columns (their \textbf{span}).\\[6pt]
      For $\begin{bsmallmatrix} 1 & 3 \\ 2 & 1 \end{bsmallmatrix}$ the columns are not parallel, so
      $C(A)$ is the \textbf{whole plane}: every $\bb$ is reachable.
    \end{column}
    \begin{column}{0.42\textwidth}
      \centering
      \begin{tikzpicture}[scale=0.6]
        \fill[burgundy!12] (-2.2,-2.2) rectangle (2.2,2.2);
        \draw[step=1, black!15, thin] (-2.2,-2.2) grid (2.2,2.2);
        \draw[->, thick, black!60] (-2.4,0)--(2.4,0);
        \draw[->, thick, black!60] (0,-2.4)--(0,2.4);
        \draw[->, very thick, burgundy] (0,0)--(1,2) node[above]{\scriptsize $\mathbf{a}_1$};
        \draw[->, very thick, royalblue] (0,0)--(2,0.667) node[right]{\scriptsize $\mathbf{a}_2$};
      \end{tikzpicture}
    \end{column}
  \end{columns}
\end{frame}

% LINE VS PLANE
\begin{frame}[t, plain]
  \burgundyheader{Line vs.\ plane --- is $\bb$ reachable?}
  \sectionlabel[goldacc]{Independence}
  \begin{block}{The rule}
    Independent (non-parallel) columns fill the whole \textbf{plane}. Parallel columns fill only a
    \textbf{line}. And $A\xx=\bb$ has a solution exactly when $\bb$ lies in $C(A)$.
  \end{block}
  For $B=\begin{bsmallmatrix} 1 & 2 \\ 2 & 4 \end{bsmallmatrix}$ column 2 is $2\times$ column 1, so
  $C(B)$ is a line: $\begin{bsmallmatrix} 3\\6 \end{bsmallmatrix}$ is reachable, but
  $\begin{bsmallmatrix} 1\\0 \end{bsmallmatrix}$ is \emph{not}.
\end{frame}

% WRAP / PREVIEW
\begin{frame}[t, plain]
  \burgundyheader{Wrap-up and what's next}
  \sectionlabel{Connections}
  \textbf{Today:} a matrix is vectors in columns; $A\xx=x_1\mathbf{a}_1+x_2\mathbf{a}_2$; the column
  space $C(A)$ is the span of the columns --- a line (parallel columns) or the whole plane
  (independent columns).\\[8pt]
  \textbf{Next --- Lesson 1.4:} build every column of a matrix from a few \emph{independent}
  columns --- the factorization $A=CR$.
\end{frame}

\end{document}
